\begin{established}[Function fields and point valuations on the cut]
\label{hyp:slice-valued-point-lift}\label{cor:generic-slice-lifting-1}
Let \(D'\) be the normalization of \(O'\) in the finite cut algebra.
Each branch of \(D'\) belongs to a factor associated to a component
\(Z\) and its valued point field contains \(\Omega=K_0(Z)\), of
transcendence degree \(s\). Equal-weight ratios are evaluated through
this inclusion and rational functions in the cutting coefficients
through the original Gauss field. Every \(f\in\Omega\) can be expressed as
\[
       f=g/P(w),\qquad \operatorname{wt}(g)=k w_L,\quad \deg P=k,
\]
where \(P\) has constant coefficients and the numerator is a cusp
section, extended by zero to the other components.
\end{established}
\begin{proof}
Choose on \(Z\) a light column \(w_b\) nonzero at its generic point.
The incidence equation solves for one coefficient in every cutting
row. Write \(X'\) for all coefficients except the pivot coefficients
\(X_{ib}\) with \(i>0\). The field on the component is
\[
 X_{ib}=-\sum_{l\ne b} X_{il}\,w_l/w_b \quad (i>0),\qquad
 K_0(\mathrm{Inc}_Z)=\Omega(X').
\]

Normalizing and
passing to a branch retains the indicated embeddings. The point-frame
calculation of Proposition~\ref{prop:selected-generic-cut} identifies
the order of each ratio with its branch order. Its homogeneous
denominator argument clears a monic equation for \(f\) over the light
image, extends \(P(w)f\) in weight \(k w_L\), and proves that its boundary
coefficients vanish. Extension by zero is defined on the disjoint
components of the smooth variety. On \(L_0\ne0\) the cut is defined
by \(J=(L_1/L_0,\ldots,L_s/L_0)\), so its weight-zero algebra is the
chart quotient by \(J\), and its weight-\(w_F\) line is free in the
frame of Lemma~\ref{hyp:slice-chart-branch-lattices}.
\end{proof}

\begin{proposition}[Lifting the completed branch lattice]
\label{prop:generic-branch-jet-lift}
For the data of Lemma~\ref{hyp:slice-valued-point-lift}, every
\(b\in B_0\) in the integral branch lattice of
Lemma~\ref{hyp:slice-chart-branch-lattices} has a
weight-\(w_F\) chart lift with integral Hasse derivatives modulo \(J^{n_0}\)
for \(n_0=\max(1,a)\). If \(s_b\) denotes this lift and \(D^{[0]}\) its restriction to the cut (the zeroth Hasse derivative), then
\[
 D^{[0]}s_b=b,\qquad n_0=\max(1,a).
\]
 The approximating
numerator may be taken at whatever degree is needed for the clearing.
\end{proposition}

\begin{proof}
\proofstep{Normalized orders in the branch field.}
At a branch write \(v\) for its integer order and \(k\) for
the residue field of \(K_0\) at \(\lambda\).
Choose a pivot \(w_b\) with least finite order among the
light columns in a common frame. All \(r_l=w_l/w_b\) on
its component have nonnegative order. Let \(n\) be the number
of variables \(X_{il}\), counting row zero. On the intersection
we have
\[
          X_{ib}=-\sum_{l\ne b} X_{il}r_l
                  \qquad(1\le i\le s).
\]
Call the other \(n-s\) variables \(X'\). All \(n\) residues
of the \(X_{il}\) are algebraically independent over \(k\)
by restriction to the original Gauss DVR. The residue field
of \(\Omega\) has transcendence degree at most \(s\) over
\(k\). Indeed any residue-independent family lifts to an
algebraically independent family over \(K_0\), by normalizing
a proposed polynomial relation at a coefficient of minimum order.

The displayed relations put all \(n\) residues in the field
generated by this residue field and the residues of \(X'\).
If \(\bar\Omega\) is the residue field of \(\Omega\), the
two dimension bounds read
\[
 n\le\operatorname{trdeg}_k\bar\Omega(\bar X')
       \le \operatorname{trdeg}_k\bar\Omega+(n-s)\le n.
\]
Thus the \(n-s\) residues in \(\bar X'\) are algebraically
independent over \(\bar\Omega\). For \(P=\sum_\alpha
a_\alpha (X')^\alpha\ne0\), divide by a coefficient \(a_\beta\)
of minimum order. Its reduction has a nonzero coefficient and
does not vanish at \(\bar X'\). If \(Q=\sum_\alpha b_\alpha(X')^\alpha\ne0\) as well, the two
normalized reductions give exactly
\[
 v(P/Q)=\min_\alpha v(a_\alpha)-\min_\alpha v(b_\alpha).
\]
In particular its values already belong to the value group of
the restricted valuation on \(\Omega\).

Evaluations of slice functions at this branch lie in the
valuation closure of \(\Omega(X')\). A homogeneous fraction
of weight zero, written with denominator \(L_0^t\), can be
evaluated by dividing its numerator and denominator by
\(w_b^t\). All ratios of the global homogeneous forms involved
then lie in \(\Omega\), while the coefficients belong to
\(K_{0,\lambda}(X)\). Any such coefficient is approximated
arbitrarily closely at the Gauss valuation by elements
of \(K_0(X)\), on approximating its finitely many polynomial
coefficients and then taking the quotient. At a finite list
of branches this approximation can be made simultaneously.
Nonvanishing denominators remain nonzero under sufficiently
close approximation.
For a denominator \(Q\ne0\) one may choose its approximation \(Q'\) so that
\[
  v(Q'-Q)>v(Q)\quad\Longrightarrow\quad v(Q')=v(Q).
\]

\proofstep{Restricting the coordinate valuation.}
In particular the restriction of \(v\) to \(\Omega\) takes
the value one. Choose an element of the slice algebra of
order one at this branch, by localization in the finite
product of fields, and approximate its value strictly closer
than its own order. Write the approximant as \(P/Q\) with coefficients \(a_\alpha,b_\beta\)
and denote the order-one element by \(z\). The ultrametric inequality
and the minimum formula give
\[
 v(z-P/Q)>v(z)=1
 \ \Longrightarrow\ 
 1=v(P/Q)=\min_\alpha v(a_\alpha)-\min_\beta v(b_\beta)
 \in v(\Omega^\times).
\] Moreover distinct branches over the
same component give inequivalent valuations on \(\Omega\).

If their restrictions were equivalent, normalize to give
the base uniformizer value one. The valuations would agree
on \(\Omega\), allowing the same pivot, and agree on every
rational expression in \(X'\) over it by the minimum rule.
For each slice element use the same homogeneous expression
at both branches and the simultaneous coefficient approximation.
Its two valuations would agree, including infinite order by
arbitrarily close approximation. Distinct branches are
distinguished by an element of their different maximal
ideals in \(D'\). This proves the assertion.

Let \(A_{\rm loc}\) be the weight-zero chart with all
denominators nonvanishing on the generic slice inverted.
Use the simultaneous branch Taylor expansion
\[
 A_{\rm loc}\longrightarrow \prod_\zeta
                           K_\zeta[[\Delta X]],
\]
where \(\Delta X\) are the increments and \(K_\zeta\) the
completed branch fields. Pull back the subrings
\(O_\zeta[[\Delta X]]\) and take the image \(C\) of this
pullback under constant-term evaluation. The Taylor
product identity makes \(C\) a subring of \(D'\).
The inverse-series formula on the Gauss ring gives
\(O'\subset C\).
A tuple of functions of the components, integral at all
branches, belongs to \(C\). Indeed write each component
function as a ratio of homogeneous forms of equal degree,
using the birational generation.

The denominator is nonzero
as a function there, hence at each of its embedded generic
points. Pad the degrees by powers of a nonzero light section
on each component, extend by zero, and sum numerators and
denominators separately.
Once all forms have the same degree, extension by zero gives the desired componentwise ratio
\[
  \left.\frac{\sum_Z g_Z}{\sum_Z h_Z}\right|_{Z'}=\frac{g_{Z'}}{h_{Z'}} .
\]
 This represents the whole tuple
by one ratio with denominator nonzero on the slice. To put
the ratio in the localized chart one first divides both forms
by the same power of \(L_0\). If \(h\) denotes this ratio of component functions, the same cancellation
holds after replacing \(X\) by \(X+\Delta X\). Thus
\[
 h(X+\Delta X)=h(X),\qquad
 [(\Delta X)^\alpha]h(X+\Delta X)=0\quad (|\alpha|>0).
\]
Its constant coefficient is the prescribed integral tuple, which is
therefore in \(C\).

\proofstep{Identifying the completed lattices.}
We show \(C=D'\). The restricted component valuations just
proved are nontrivial and pairwise inequivalent on each
component. Weak approximation at this finite set realizes
arbitrary assigned integer orders, using an element of order
one locally. It can also approximate a prescribed value
at one branch while prescribing integrality at all the other
branches. Do this component by component and use the homogeneous
ratio just constructed. Thus \(C\) realizes all nonnegative
order profiles, and a locally integral component function has
the same residue at a given branch as some element of \(C\).
\cite{weak-approximation}

Fix a branch \(\zeta\). For an element of \(D'\) approximate
its value there by \(P/Q\in\Omega(X')\) with positive valuation
of the difference. After dividing both polynomials by one
minimum-order coefficient of \(Q\), the Gauss formula gives

\[
  v_\zeta(Q)=0,
\]
 \(v_\zeta(P)\ge0\). Thus their coefficients
are integral at this branch. Replace their residues by the
elements of \(C\) supplied by weak approximation; the variables
already belong to \(O'\). Write \(C_\zeta\) for the image of \(C\) in the branch residue
field. After coefficient replacement the desired residue is
\[
            \bar P/\bar Q\in\operatorname{Frac}(C_\zeta),
                      \qquad \bar P,\bar Q\in C_\zeta,\quad \bar Q\ne0 .
\]
The field of \(O'\)-residues lies in \(C_\zeta\). Since \(C_\zeta\)
is a domain finite-dimensional over that field, multiplication by
each nonzero element is surjective, so \(C_\zeta\) is a field.
Thus the reduction of \(C\) surjects onto the residue field at
each branch.

\proofstep{Separating the branch idempotents.}
For completeness write \(\mathfrak q_\zeta\) for the branch
prime in \(D'\). The ideals \(C\cap\mathfrak q_\zeta\) are
maximal by residue surjectivity and distinct by the order
profiles. At each branch take an element of \(C\) of order
one. Lift the residue digits successively with its powers
as place values. This surjects onto the local quotient modulo
\(\mathfrak q_\zeta^{e_\zeta}\). Put \(\mathfrak m_\zeta=C\cap\mathfrak q_\zeta\).
Distinctness of these maximal ideals permits simultaneous
choice of the lifted digits by the Chinese remainder theorem.
Since \(\mathfrak m_\zeta^{e_\zeta}\) maps to zero modulo
\(\mathfrak q_\zeta^{e_\zeta}\), the maps give a surjection
\[
 C/\!\prod_\zeta\mathfrak m_\zeta^{e_\zeta}
   \longrightarrow
    \prod_\zeta D'/\mathfrak q_\zeta^{e_\zeta}
       \simeq D'/\pi D'.
\]
The last equality uses the branch ramification exponents.
Thus the finite \(O'\)-module \(D'/C\) has zero reduction
modulo \(\pi\), and Nakayama gives \(C=D'\).

\proofstep{Clearing the common denominator.}
Finally choose a ratio of homogeneous forms, independent
of \(X\), of order \(-c_\zeta\) at each branch by weak
approximation. Multiplying \(g_0\) by this ratio gives a
frame of point order zero everywhere and with vanishing
positive Taylor coefficients. In this frame the shifted ball
\(B_0\) is just \(D'\). The equality \(C=D'\) supplies a
localized section lifting every \(b\in B_0\) with integral
Hasse readings.

To return to the chart before localization
take a common weight-zero denominator \(d\) invertible on
the finite reduced slice. Lift its inverse modulo \(J\) to
a chart function \(d_1\). Then
\[
                 d_1\sum_{i=0}^{n_0-1}(1-dd_1)^i
              \equiv d^{-1}\pmod {J^{n_0}},\qquad
                 n_0=\max(1,a).
\]
After clearing the possible equal-weight denominator of the
frame in the same manner, this gives an actual weight-\(w_F\)
section modulo \(J^{n_0}\). Its restriction is \(b\); Hasse
derivatives of total degree less than \(n_0\) of the difference
vanish on the slice by the product formula. This proves the
lifting requirement at every order for the stated geometric
realization. In particular no effective upper bound on the
degree of the approximating chart numerator is needed in
the pole-length calculation. \label{cor:generic-slice-lifting-2}
\end{proof}
